\begin{definition}[Terminal spectral refinement of a fusion decomposition]%
    \label{def:mpdo_fusion_clause_terminal_spectral_refinement}
    \lean{MPOTensor.BNTFusionTensorClause.terminalMatrix}
    \lean{MPOTensor.BNTFusionTensorClause.TerminalSpectralIndex}
    \lean{MPOTensor.BNTFusionTensorClause.terminalEigenvalue}
    \lean{MPOTensor.BNTFusionTensorClause.terminalEigenProjection}
    \lean{MPOTensor.BNTFusionTensorClause.terminalEigenProjectionFamily}
    \lean{MPOTensor.BNTFusionTensorClause.%
        recursiveTerminalEigenProjectorQ}
    \lean{MPOTensor.BNTFusionTensorClause.%
        topologicalSpectralProjectorSucc}
    \lean{MPOTensor.BNTFusionTensorClause.%
        HasTerminalSpectralProjectorRefinement}
    \leanok
    \uses{def:mpdo, def:mpdo_bnt_fusion_tensor_clause, def:mpdo_topological_density_operator, def:mpdo_topological_projector_recursion, thm:mpdo_terminal_transfer_positivity}
    Let $M$ generate MPDOs, and fix a fusion decomposition of its vertical
    BNT $\{M_\alpha\}_\alpha$ with positive diagonal matrices
    $\chi_{\alpha,\beta,\gamma}$. For every terminal label $\gamma$, close the
    horizontal operator leg of $M_\gamma$ to obtain a positive semidefinite
    bond matrix $K_\gamma$. Choose an orthonormal eigenbasis of each
    $K_\gamma$. Its spectral indices are pairs $s=(\gamma,k)$; write
    $\lambda_s$ for the corresponding eigenvalue and $E_s$ for the rank-one
    orthogonal projection onto the chosen eigenvector. Recursively transport
    this family through every fusion history and the adjoint sequential
    coisometry. Denote the resulting operator, assembled over all labels and
    multiplicities at positive length $N$, by $P_s^{(N)}$.

    Only nonzero product sectors occur in the fusion decomposition. A normal
    label absent from a fixed product pair has zero fusion multiplicity. See
    \cite{gap:cpsv16_bnt_uniqueness_zero_coefficient}.

    The fixed-pair decomposition in Figure~11 requires a support correction:
    a fixed product pair may have no nonzero sectors. See
    \cite{gap:cpsv16_figure11_per_pair_support}.

    The fusion map restricted to the retained rows is a coisometry onto the
    direct sum of nonzero sectors, and its adjoint gives exact reconstruction.
    This corrects the fusion-map convention in Figure~11. See
    \cite{gap:cpsv16_figure11_fusion_coisometry}.
\end{definition}

\begin{definition}[Terminal spectral refinement for an RFP MPDO]%
    \label{def:mpdo_terminal_spectral_refinement}
    \lean{MPOTensor.rfpBNTFusionTensorClause}
    \leanok
    \uses{def:horizontal_cf_mpdo, def:mpdo, def:is_rfp_via_ts, def:mpdo_bnt_fusion_tensor_clause, def:mpdo_fusion_clause_terminal_spectral_refinement, thm:mpdo_horizontal_rfp_to_bnt_fusion_tensor}
    For an MPDO in normalized BNT-refined horizontal form satisfying the
    RFP condition, fix the fusion decomposition supplied by that
    construction. For every terminal label $\gamma$, close the horizontal
    operator leg of $M_\gamma$ to obtain a matrix $K_\gamma$ on its bond
    space. Choose an orthonormal eigenbasis of each $K_\gamma$, with spectral
    indices $s=(\gamma,k)$ satisfying $0\leq k<D_\gamma$.
    Write $\lambda_s$ for the corresponding eigenvalue and $E_s$ for the
    rank-one orthogonal projection onto the chosen eigenvector, extended by
    zero on terminal labels different from~$\gamma$. Recursively transport
    this terminal family through every fusion history, then through the
    adjoint sequential fusion coisometry, and take the direct sum over all
    sitewise label and multiplicity configurations. Denote the resulting
    operator at positive chain length $N$ by $P_s^{(N)}$.

    The normalized BNT-refined horizontal hypothesis is stronger than CF
    as defined by CPSV.
    % Scope tracking: docs/paper-gaps/cpgsv17_vertical_cf_grouping.tex.
\end{definition}

\begin{theorem}[Terminal spectral projectors from a fusion decomposition]%
    \label{thm:mpdo_fusion_clause_terminal_spectral_projector_refinement}
    \lean{MPOTensor.BNTFusionTensorClause.terminalEigenvalue_nonneg}
    \lean{MPOTensor.BNTFusionTensorClause.%
        topologicalSpectralProjectorSucc_isOrthogonalProjection}
    \lean{MPOTensor.BNTFusionTensorClause.%
        topologicalSpectralProjectorSucc_mul_eq_zero}
    \lean{MPOTensor.BNTFusionTensorClause.%
        sum_terminalEigenvalue_smul_topologicalSpectralProjectorSucc}
    \lean{MPOTensor.BNTFusionTensorClause.%
        topologicalMultiplicityWeightFactorSucc_commutes_spectralProjector}
    \lean{MPOTensor.BNTFusionTensorClause.%
        topologicalDensityOperatorSucc_eq_sum_terminalSpectralProjector}
    \lean{MPOTensor.BNTFusionTensorClause.%
        hasTerminalSpectralProjectorRefinement}
    \leanok
    \uses{def:mpdo, def:mpdo_bnt_fusion_tensor_clause, def:mpdo_fusion_clause_terminal_spectral_refinement, thm:mpdo_fusion_clause_topological_density_decomposition, cor:mpdo_fusion_clause_topological_density_factor_commutator}
    Let $M$ generate MPDOs, fix a fusion decomposition on the nonzero product
    sectors, and suppose that its structure coefficients
    $c_{\alpha,\beta,\gamma}^{(L)}
        =\tr(\chi_{\alpha,\beta,\gamma}^{L})$ are independent of $L$ for all
    positive lengths, with $\chi_{\alpha,\beta,\gamma}$ positive diagonal.
    Then every terminal eigenvalue $\lambda_s$ is non-negative. For every
    $N>0$, the operators $P_s^{(N)}$ are pairwise orthogonal projections and
    \begin{align}
        T_N          & =\sum_s\lambda_sP_s^{(N)},
        \notag                                       \\
        A_NP_s^{(N)} & =P_s^{(N)}A_N,
        \notag                                       \\
        R_N          & =\sum_s\lambda_sA_NP_s^{(N)}.
        \notag
    \end{align}
    This is the terminal spectral decomposition of
    \cite[lines~1003--1012]{Cirac2016MPDO_arXiv}; it makes no Hamiltonian or
    Gibbs-state assertion.

    Only nonzero product sectors occur in the fusion decomposition. A normal
    label absent from a fixed product pair has zero fusion multiplicity. See
    \cite{gap:cpsv16_bnt_uniqueness_zero_coefficient}.

    The fixed-pair decomposition in Figure~11 requires a support correction:
    a fixed product pair may have no nonzero sectors. See
    \cite{gap:cpsv16_figure11_per_pair_support}.

    The fusion map restricted to the retained rows is a coisometry onto the
    direct sum of nonzero sectors, and its adjoint gives exact reconstruction.
    This corrects the fusion-map convention in Figure~11. See
    \cite{gap:cpsv16_figure11_fusion_coisometry}.
\end{theorem}
\begin{proof}\leanok
    Positivity of the physical density operator makes each terminal bond
    matrix positive semidefinite. Its spectral theorem supplies the
    non-negative eigenvalues and orthogonal rank-one projections. Length
    independence makes every retained diagonal entry of
    $\chi_{\alpha,\beta,\gamma}$ equal to one, so recursive coisometric
    transport preserves orthogonality and the spectral sum. The
    density-factor commutator gives the final two identities.
\end{proof}

\begin{theorem}[Terminal spectral projectors for an RFP MPDO with
        length-independent structure coefficients]%
    \label{thm:mpdo_terminal_spectral_projector_refinement}
    \lean{MPOTensor.terminalSpectralProjectorRefinement_of_isRFPViaTS}
    \leanok
    \uses{def:horizontal_cf_mpdo, def:mpdo, def:is_rfp_via_ts, def:mpdo_terminal_spectral_refinement, thm:mpdo_horizontal_rfp_to_bnt_fusion_tensor, thm:mpdo_fusion_clause_terminal_spectral_projector_refinement}
    Let $M$ generate MPDOs, be in normalized BNT-refined horizontal form,
    and satisfy the RFP condition. Fix the fusion decomposition supplied by
    that construction, and suppose that its structure coefficients
    $c_{\alpha,\beta,\gamma}^{(L)}
        =\tr(\chi_{\alpha,\beta,\gamma}^{L})$, with positive diagonal
    $\chi_{\alpha,\beta,\gamma}$, are independent of $L$ for all positive
    lengths.
    Then $K_\gamma$ is positive semidefinite for every terminal label, so
    $\lambda_s\geq0$. Write the positive chain length as $N=n+1$,
    where $n\in\N$. For every $N>0$,
    \begin{align}
        (P_s^{(N)})^\dagger
         & =P_s^{(N)},
        \notag         \\
        (P_s^{(N)})^2
         & =P_s^{(N)}.
        \notag
    \end{align}
    For distinct spectral indices $s$ and $t$, $P_s^{(N)}P_t^{(N)}=0$.
    The recursive factor and multiplicity weights satisfy
    \begin{align}
        T_N
         & =\sum_s\lambda_s P_s^{(N)},
        \label{eq:rfp_terminal_spectral_sum} \\
        A_NP_s^{(N)}
         & =P_s^{(N)}A_N.
        \label{eq:rfp_spectral_commutation}
    \end{align}
    Consequently,
    \begin{align}
        R_N
         & =\sum_s\lambda_s A_NP_s^{(N)}.
        \label{eq:rfp_density_spectral_sum}
    \end{align}
    Length independence is assumed for the coefficients of the fixed fusion
    decomposition supplied by the RFP construction, not for an arbitrary
    fusion decomposition.

    The normalized BNT-refined horizontal hypothesis is stronger than CF
    as defined by CPSV.
    % Scope tracking: docs/paper-gaps/cpgsv17_vertical_cf_grouping.tex.

    This theorem gives the spectral decomposition in
    \cite[lines~1003--1012]{Cirac2016MPDO_arXiv}; it does not assert the
    commuting nearest-neighbor Hamiltonian or Gibbs-state conclusion beginning
    at line~1013.
\end{theorem}
\begin{proof}\leanok
    \uses{thm:mpdo_topological_projector_recursion, thm:mpdo_topological_density_decomposition, cor:mpdo_topological_density_factor_commutator}
    Positivity of the one-site density operator passes to every fixed BNT
    block; division by its strictly positive multiplicity weight shows that
    each $K_\gamma$ is positive semidefinite. Apply the finite-dimensional
    spectral theorem to obtain the non-negative eigenvalues and rank-one
    orthogonal projections. Length independence makes every retained diagonal
    entry of $\chi_{\alpha,\beta,\gamma}$ equal to one, so the recursive
    direct sum preserves idempotence and pairwise orthogonality. Write
    $Q_{N,c,s}$ for the recursive spectral block in copy
    configuration~$c$. The coisometry identity
    $W_{N,c}W_{N,c}^\dagger=\Id$ gives
    \begin{align}
        (W_{N,c}^\dagger Q_{N,c,s}W_{N,c})^\dagger
         & =W_{N,c}^\dagger Q_{N,c,s}W_{N,c},
        \notag                                \\
        (W_{N,c}^\dagger Q_{N,c,s}W_{N,c})^2
         & =W_{N,c}^\dagger Q_{N,c,s}W_{N,c}.
        \notag
    \end{align}
    Thus transport through the adjoint sequential coisometry preserves the
    projector identities. Linearity gives
    (\ref{eq:rfp_terminal_spectral_sum}). Finally, $A_N$ is scalar on every
    copy-configuration block, hence it commutes with each $P_s^{(N)}$ and
    gives (\ref{eq:rfp_density_spectral_sum}).
\end{proof}

\begin{definition}[Commuting Gibbs form on the retained vertical space for
        chains of at least two sites]%
    \label{def:mpdo_topological_gibbs_data_at_least_two}
    \lean{MPOTensor.BNTFusionTensorClause.%
        terminalSpectralIndex_card_le_physicalDim}
    \lean{MPOTensor.BNTFusionTensorClause.terminalSpectralEmbedding}
    \lean{MPOTensor.BNTFusionTensorClause.%
        physicalIndexedTerminalEigenvalue}
    \lean{MPOTensor.BNTFusionTensorClause.%
        physicalIndexedTopologicalSpectralProjectorSucc}
    \lean{MPOTensor.BNTFusionTensorClause.retainedMultiplicityWeightEntry}
    \lean{MPOTensor.BNTFusionTensorClause.retainedMultiplicityEnergyEntry}
    \lean{MPOTensor.BNTFusionTensorClause.topologicalGibbsLocalTerm}
    \lean{MPOTensor.BNTFusionTensorClause.%
        topologicalGibbsHamiltonianSuccSucc}
    \lean{MPOTensor.BNTFusionTensorClause.HasTopologicalGibbsDecomposition}
    \leanok
    \uses{def:mpdo_fusion_clause_terminal_spectral_refinement}
    Let $M$ generate MPDOs, fix a fusion decomposition on the nonzero product
    sectors, and let $\mu$ be its strictly positive diagonal multiplicity
    matrix on the retained one-site space. Define
    $h=(-\log\mu)\otimes\Id$ on two adjacent sites, where $\log\mu$ is the
    diagonal matrix obtained by taking the logarithm of each diagonal entry.
    For a chain of length $L=N+2$, with $N\in\N$, let
    $H_L=\sum_{i=0}^{L-1}h_{i,i+1}$, with periodic indices.

    The terminal spectral family has cardinality at most the physical
    one-site dimension $d$. Fix an injection of its indices into
    $\{0,\ldots,d-1\}$ and extend the eigenvalues and projectors by zero
    outside the image.

    The Hamiltonian and projectors in this definition act on the retained
    nonzero vertical sectors. Applying the adjoint of the coisometry
    restricted to the retained rows reconstructs the physical density
    operator; this reconstruction alone does not specify a Hamiltonian or
    projectors on the original physical space. The distinction is recorded in
    \cite{gap:cpsv16_topological_gibbs_physical_complement}.

    The local interaction is defined on two spins, so this definition applies
    to lengths $L=N+2$. The length-one boundary is recorded in
    \cite{gap:cpsv16_topological_gibbs_length_one}.
\end{definition}

\begin{theorem}[Fusion decomposition gives a commuting Gibbs decomposition
        on the retained vertical space for chains of at least two sites]%
    \label{thm:mpdo_fusion_clause_topological_gibbs_at_least_two}
    \lean{MPOTensor.BNTFusionTensorClause.%
        topologicalGibbsLocalTerm_isHermitian}
    \lean{MPOTensor.BNTFusionTensorClause.%
        topologicalGibbsBondSuccSucc_commute}
    \lean{MPOTensor.BNTFusionTensorClause.%
        exp_neg_topologicalGibbsHamiltonianSuccSucc}
    \lean{MPOTensor.BNTFusionTensorClause.%
        physicalIndexedTopologicalSpectralProjectorSucc_isOrthogonalProjection}
    \lean{MPOTensor.BNTFusionTensorClause.%
        physicalIndexedTopologicalSpectralProjectorSucc_commutes_gibbsFactor}
    \lean{MPOTensor.BNTFusionTensorClause.%
        topologicalDensityOperatorSucc_eq_sum_physicalIndexedProjector_mul_gibbsFactor}
    \lean{MPOTensor.BNTFusionTensorClause.%
        singleKrausMap_physicalIndexedGibbsDecomposition_eq_mpo}
    \lean{MPOTensor.BNTFusionTensorClause.hasTopologicalGibbsDecomposition}
    \leanok
    \uses{def:mpdo, def:mpdo_bnt_fusion_tensor_clause, def:mpdo_topological_gibbs_data_at_least_two, thm:mpdo_fusion_clause_terminal_spectral_projector_refinement}
    Let $M$ generate MPDOs, and fix a fusion decomposition on the nonzero
    product sectors of a vertical canonical decomposition of $M$.
    Suppose that the structure coefficients
    $c_{\alpha,\beta,\gamma}^{(L)}
        =\tr(\chi_{\alpha,\beta,\gamma}^{L})$, with positive diagonal
    $\chi_{\alpha,\beta,\gamma}$, are independent of $L$ for all positive
    lengths.
    There are non-negative numbers $\lambda_0,\ldots,\lambda_{d-1}$,
    independent of the chain length. The two-site matrix $h$ is Hermitian,
    and for every $L\geq2$ its periodic translates commute pairwise and satisfy
    \begin{align}
        H_L
         & =\sum_{j=0}^{L-1}h_{j,j+1},
        \label{eq:rfp_gibbs_hamiltonian} \\
        e^{-H_L}
         & =\mu^{\otimes L}.
        \label{eq:rfp_gibbs_factor}
    \end{align}
    There are orthogonal projections
    $P_0^{(L)},\ldots,P_{d-1}^{(L)}$ such that
    \begin{align}
        [P_i^{(L)},e^{-H_L}]
         & =0,
        \label{eq:rfp_gibbs_projector_commutation}      \\
        R_L
         & =\sum_{i=0}^{d-1}\lambda_iP_i^{(L)}e^{-H_L}.
        \label{eq:rfp_retained_gibbs_sum}
    \end{align}
    Applying the adjoint of the coisometry restricted to the retained rows
    at every site reconstructs $\rho^{(L)}(M)$ from the right-hand side of
    (\ref{eq:rfp_retained_gibbs_sum}).

    Only nonzero product sectors occur in the fusion decomposition. A normal
    label absent from a fixed product pair has zero fusion multiplicity. See
    \cite{gap:cpsv16_bnt_uniqueness_zero_coefficient}.

    The fixed-pair decomposition in Figure~11 requires a support correction:
    a fixed product pair may have no nonzero sectors, and no unsupported
    normal sector is inserted. See
    \cite{gap:cpsv16_figure11_per_pair_support}.

    The fusion map restricted to the retained rows is a coisometry onto the
    direct sum of nonzero sectors, and its adjoint gives exact reconstruction
    of the product tensor. This corrects the fusion-map convention in
    Figure~11. See
    \cite{gap:cpsv16_figure11_fusion_coisometry}.

    The formula is an equality for $R_L$ on the retained vertical space, not
    the physical-space equality asserted in
    \cite[lines~1013--1016]{Cirac2016MPDO_arXiv}. Exact reconstruction by the
    adjoint coisometry alone does not establish the existence of physical
    projectors and a finite physical Hamiltonian. See
    \cite{gap:cpsv16_topological_gibbs_physical_complement}.

    The conclusion holds for $L=N+2$, with $N\in\N$. The definition of the
    local interaction in \cite{Cirac2016MPDO_arXiv} uses two spins but does
    not state a length-one convention; see
    \cite{gap:cpsv16_topological_gibbs_length_one}.
\end{theorem}
\begin{proof}\leanok
    \uses{thm:mpdo_fusion_clause_topological_density_decomposition, cor:mpdo_fusion_clause_topological_density_factor_commutator, thm:mpdo_fusion_clause_terminal_spectral_projector_refinement}
    Each diagonal entry of $\mu$ is strictly positive, so its negative
    logarithm is real. The translated terms are diagonal, and each site
    contributes once to their sum. Exponentiating gives
    (\ref{eq:rfp_gibbs_factor}).

    Positivity of the multiplicities embeds one terminal mode for every
    terminal spectral index into the retained coordinates. The rank of the
    coisometry restricted to the retained rows is at most $d$, so the terminal
    family has at most $d$ members. Extend the eigenvalues and spectral
    projectors by zero along an injection into $\{0,\ldots,d-1\}$. The added
    zero projectors preserve self-adjointness, idempotence, commutation, and
    the spectral sum. Combine that sum with (\ref{eq:rfp_gibbs_factor}) and
    the exact adjoint reconstruction of the retained density operator.
\end{proof}

\begin{theorem}[Commuting Gibbs decomposition on the retained vertical space
        for chains of at least two sites]%
    \label{thm:mpdo_topological_gibbs_at_least_two}
    \lean{MPOTensor.topologicalGibbsDecomposition_of_isRFPViaTS}
    \leanok
    \uses{def:horizontal_cf_mpdo, def:mpdo, def:is_rfp_via_ts, thm:mpdo_horizontal_rfp_to_bnt_fusion_tensor, thm:mpdo_fusion_clause_topological_gibbs_at_least_two}
    Let $M$ generate MPDOs, be in normalized BNT-refined horizontal form,
    and satisfy the RFP condition. Select the fusion decomposition on the
    nonzero product sectors supplied by that condition, and suppose that the
    structure coefficients $c_{\alpha,\beta,\gamma}^{(L)}
        =\tr(\chi_{\alpha,\beta,\gamma}^{L})$ of this same decomposition are
    independent of $L$ for all positive lengths. Then the commuting Gibbs
    decomposition on the retained vertical space in
    Theorem~\ref{thm:mpdo_fusion_clause_topological_gibbs_at_least_two}
    holds for every $L=N+2$, with $N\in\N$.

    The normalized BNT-refined horizontal hypothesis is stronger than CF
    as defined by CPSV.

    The Hamiltonian and projectors act on the retained nonzero vertical
    sectors. Applying the adjoint of the coisometry restricted to the
    retained rows reconstructs the physical density operator; this
    reconstruction alone does not specify a Hamiltonian or projectors on the
    original physical space. See
    \cite{gap:cpsv16_topological_gibbs_physical_complement}.

    The conclusion holds for $L=N+2$; no length-one convention for the
    two-site interaction is asserted. See
    \cite{gap:cpsv16_topological_gibbs_length_one}.
\end{theorem}
\begin{proof}\leanok
    Select the fusion decomposition on the nonzero product sectors supplied
    by the RFP condition and apply
    Theorem~\ref{thm:mpdo_fusion_clause_topological_gibbs_at_least_two}
    with the assumed length independence of its structure coefficients.
\end{proof}

\begin{theorem}[Fusion decomposition gives a physical commuting Gibbs
        decomposition for chains of at least two sites]%
    \label{thm:mpdo_fusion_clause_physical_coordinate_gibbs_at_least_two}
    \lean{MPOTensor.BNTFusionTensorClause.%
        hasPhysicalTopologicalGibbsDecomposition}
    \leanok
    \uses{def:mpdo, def:mpdo_bnt_fusion_tensor_clause, def:mpdo_topological_gibbs_data_at_least_two, thm:mpdo_fusion_clause_topological_gibbs_at_least_two}
    Let $M$ generate MPDOs, and fix a fusion decomposition on the nonzero
    product sectors of a vertical canonical decomposition of $M$.
    Suppose that the structure coefficients
    $c_{\alpha,\beta,\gamma}^{(L)}
        =\tr(\chi_{\alpha,\beta,\gamma}^{L})$, with positive diagonal
    $\chi_{\alpha,\beta,\gamma}$, are independent of $L$ for all positive
    lengths. There are non-negative numbers $\lambda_0,\ldots,\lambda_{d-1}$
    and one fixed Hermitian two-site matrix $\widehat h$, all independent
    of the chain length. For every $L\geq2$, let
    $\widehat H_L=\sum_{j=0}^{L-1}\widehat h_{j,j+1}$, with periodic indices.
    The translates of $\widehat h$ commute pairwise, and there are orthogonal
    projections $\widehat P_0^{(L)},\ldots,\widehat P_{d-1}^{(L)}$ on the
    physical chain such that
    \begin{align}
        [\widehat h_{j-1,j},\widehat h_{j,j+1}]
         & =0,
        \notag                                          \\
        [\widehat P_i^{(L)},e^{-\widehat H_L}]
         & =0,
        \notag                                          \\
        \rho^{(L)}(M)
         & =\sum_{i=0}^{d-1}\lambda_i\widehat P_i^{(L)}
        e^{-\widehat H_L}.
        \label{eq:rfp_physical_gibbs_sum}
    \end{align}
    This is the implication from the fusion decomposition in the proof of
    \cite[lines~999--1016]{Cirac2016MPDO_arXiv}.

    Only nonzero product sectors occur in the fusion decomposition. A normal
    label absent from a fixed product pair has zero fusion multiplicity. See
    \cite{gap:cpsv16_bnt_uniqueness_zero_coefficient}.

    The fixed-pair decomposition in Figure~11 requires a support correction:
    a fixed product pair may have no nonzero sectors, and no unsupported
    normal sector is inserted. See
    \cite{gap:cpsv16_figure11_per_pair_support}.

    The fusion map restricted to the retained rows is a coisometry onto the
    direct sum of nonzero sectors, and its adjoint gives exact reconstruction
    of the product tensor. This corrects the fusion-map convention in
    Figure~11. See
    \cite{gap:cpsv16_figure11_fusion_coisometry}.

    To define the operators on the full physical space, transport the retained
    one-site energy through the adjoint coisometry and extend it by zero
    energy on the orthogonal complement. Transport the retained projectors
    through the adjoint sitewise coisometry. This supplies the physical
    complement correction described in
    \cite{gap:cpsv16_topological_gibbs_physical_complement}.

    The conclusion holds for $L=N+2$, with $N\in\N$. The definition of the
    local interaction in \cite{Cirac2016MPDO_arXiv} uses two spins but gives
    no length-one convention. See
    \cite{gap:cpsv16_topological_gibbs_length_one}.
\end{theorem}
\begin{proof}\leanok
    Let $V$ be the coisometry restricted to the retained rows and let $k$ be
    the retained logarithmic multiplicity energy. The physical energy
    $K=V^\dagger kV$ is Hermitian and
    \begin{align}
        e^{-K}
         & =V^\dagger e^{-k}V+(\Id-V^\dagger V).
        \label{eq:rfp_physical_weight}
    \end{align}
    Taking the fixed two-site term $K\otimes\Id$, its periodic translates act
    on separate one-site factors, hence commute, and their exponential is the
    tensor power of (\ref{eq:rfp_physical_weight}).

    For the sitewise coisometry $W=V^{\otimes L}$, transport each retained
    projector by $\widehat P_i^{(L)}=W^\dagger P_i^{(L)}W$. The identity
    $WW^\dagger=\Id$ preserves self-adjointness, idempotence, and pairwise
    orthogonality. The one-site intertwining identities transport the
    retained commutator to the physical Gibbs factor. Applying the same
    identities term by term to (\ref{eq:rfp_retained_gibbs_sum}), and then
    using exact adjoint reconstruction, gives
    (\ref{eq:rfp_physical_gibbs_sum}).
\end{proof}

\begin{theorem}[Physical commuting Gibbs decomposition for chains of at least
        two sites]%
    \label{thm:mpdo_physical_coordinate_gibbs_at_least_two}
    \lean{MPOTensor.physicalTopologicalGibbsDecomposition_of_isRFPViaTS}
    \leanok
    \uses{def:horizontal_cf_mpdo, def:mpdo, def:is_rfp_via_ts, thm:mpdo_horizontal_rfp_to_bnt_fusion_tensor, thm:mpdo_fusion_clause_physical_coordinate_gibbs_at_least_two}
    Let $M$ generate MPDOs, be in normalized BNT-refined horizontal form,
    and satisfy the RFP condition. Select the fusion decomposition on the
    nonzero product sectors supplied by that condition, and suppose that the
    structure coefficients $c_{\alpha,\beta,\gamma}^{(L)}
        =\tr(\chi_{\alpha,\beta,\gamma}^{L})$ of this same decomposition are
    independent of $L$ for all positive lengths. Then the physical commuting
    Gibbs decomposition of
    Theorem~\ref{thm:mpdo_fusion_clause_physical_coordinate_gibbs_at_least_two}
    holds for every $L=N+2$, with $N\in\N$.

    The normalized BNT-refined horizontal hypothesis is stronger than CF
    as defined by CPSV.

    The retained one-site energy is transported through the adjoint
    coisometry and extended by zero energy on the orthogonal complement.
    The retained projectors are transported through the adjoint sitewise
    coisometry. This supplies the physical complement correction described in
    \cite{gap:cpsv16_topological_gibbs_physical_complement}.

    The conclusion holds for $L=N+2$; no length-one convention for the
    two-site interaction is asserted. See
    \cite{gap:cpsv16_topological_gibbs_length_one}.
\end{theorem}
\begin{proof}\leanok
    Select the fusion decomposition on the nonzero product sectors supplied
    by the RFP condition and apply
    Theorem~\ref{thm:mpdo_fusion_clause_physical_coordinate_gibbs_at_least_two}
    with the assumed length independence of its structure coefficients.
\end{proof}

\begin{definition}[Fusion decomposition for an RFP tensor in CPSV canonical
        form]%
    \label{def:mpdo_literal_cpsv_rfp_bnt_fusion_tensor_clause}
    \lean{MPOTensor.cpsvRFPBNTFusionTensorClause}
    \leanok
    \uses{def:mpdo, def:cpsv_canonical_form, def:is_rfp_via_ts, def:mpdo_bnt_fusion_tensor_clause, thm:mpdo_literal_cpsv_rfp_to_bnt_fusion_tensor}
    Let $M$ generate MPDOs, suppose that its MPS tensor with paired physical
    indices is in CF as defined by CPSV, and suppose that $M$ satisfies the
    RFP condition. Choose one fusion decomposition of the vertical BNT on
    the nonzero product sectors furnished by these hypotheses and denote it
    by $\mathcal F_{\mathrm{CPSV}}(M)$.

    Only nonzero product sectors occur in the selected fusion decomposition.
    A normal label absent from a fixed product pair has zero fusion
    multiplicity. See
    \cite{gap:cpsv16_bnt_uniqueness_zero_coefficient}.

    The fixed-pair decomposition in Figure~11 requires a support correction:
    a fixed product pair may have no nonzero sectors, and no unsupported
    normal sector is inserted. See
    \cite{gap:cpsv16_figure11_per_pair_support}.

    The fusion map restricted to the retained rows is a coisometry onto the
    direct sum of nonzero sectors, and its adjoint gives exact reconstruction
    of the product tensor. This corrects the fusion-map convention in
    Figure~11. See
    \cite{gap:cpsv16_figure11_fusion_coisometry}.
\end{definition}

\begin{theorem}[Physical commuting Gibbs decomposition in CPSV canonical form
        for chains of at least two sites]%
    \label{thm:mpdo_literal_cpsv_physical_coordinate_gibbs_at_least_two}
    \lean{MPOTensor.%
        physicalTopologicalGibbsDecomposition_of_isRFPViaTS_of_cpsvCanonicalForm}
    \leanok
    \uses{def:mpdo, def:cpsv_canonical_form, def:is_rfp_via_ts, def:mpdo_bnt_fusion_tensor_clause, def:mpdo_literal_cpsv_rfp_bnt_fusion_tensor_clause, def:mpdo_topological_gibbs_data_at_least_two, thm:mpdo_literal_cpsv_rfp_to_bnt_fusion_tensor, thm:mpdo_fusion_clause_physical_coordinate_gibbs_at_least_two}
    Let $M$ generate MPDOs, suppose that its MPS tensor with paired physical
    indices is in CF as defined by CPSV, and suppose that $M$ satisfies the
    RFP condition. Select the fusion decomposition on the nonzero product
    sectors supplied by these three hypotheses, and assume that the structure
    coefficients $c_{\alpha,\beta,\gamma}^{(L)}
        =\tr(\chi_{\alpha,\beta,\gamma}^{L})$ of this same decomposition are
    independent of $L$ for all positive lengths.

    There are non-negative numbers
    $\lambda_0,\ldots,\lambda_{d-1}$ and one fixed Hermitian two-site matrix
    $\widehat h$, all independent of the chain length. For every $L=N+2$,
    with $N\in\N$, let
    $\widehat H_L=\sum_{j=0}^{L-1}\widehat h_{j,j+1}$, with periodic indices.
    The translates of $\widehat h$ commute pairwise, and there are pairwise
    orthogonal projections
    $\widehat P_0^{(L)},\ldots,\widehat P_{d-1}^{(L)}$ on the physical chain
    such that
    \begin{align}
        [\widehat h_{j-1,j},\widehat h_{j,j+1}]
         & =0,
        \notag                                          \\
        [\widehat P_i^{(L)},e^{-\widehat H_L}]
         & =0,
        \notag                                          \\
        \rho^{(L)}(M)
         & =\sum_{i=0}^{d-1}\lambda_i\widehat P_i^{(L)}
        e^{-\widehat H_L}.
        \notag
    \end{align}
    This is the physical-space conclusion of
    \cite[lines~1013--1016]{Cirac2016MPDO_arXiv}, under the CF and RFP
    hypotheses of
    \cite[Theorem~4.14, lines~972--993]{Cirac2016MPDO_arXiv}.

    Only nonzero product sectors occur in the selected fusion decomposition.
    A normal label absent from a fixed product pair has zero fusion
    multiplicity. See
    \cite{gap:cpsv16_bnt_uniqueness_zero_coefficient}.

    The fixed-pair decomposition in Figure~11 requires a support correction:
    a fixed product pair may have no nonzero sectors, and no unsupported
    normal sector is inserted. See
    \cite{gap:cpsv16_figure11_per_pair_support}.

    The fusion map restricted to the retained rows is a coisometry onto the
    direct sum of nonzero sectors, and its adjoint gives exact reconstruction
    of the product tensor. This corrects the fusion-map convention in
    Figure~11. See
    \cite{gap:cpsv16_figure11_fusion_coisometry}.

    The retained one-site energy is transported through the adjoint
    coisometry and extended by zero energy on the orthogonal complement.
    The retained projectors are transported through the adjoint sitewise
    coisometry. This supplies the physical complement correction described in
    \cite{gap:cpsv16_topological_gibbs_physical_complement}.

    The conclusion holds for $L=N+2$. The definition of the local interaction
    in \cite{Cirac2016MPDO_arXiv} uses two spins but gives no length-one
    convention. See
    \cite{gap:cpsv16_topological_gibbs_length_one}.
\end{theorem}
\begin{proof}\leanok
    Choose the fusion decomposition on the nonzero product sectors supplied by
    Theorem~\ref{thm:mpdo_literal_cpsv_rfp_to_bnt_fusion_tensor}. Its
    structure coefficients are length independent by hypothesis. Apply
    Theorem~\ref{thm:mpdo_fusion_clause_physical_coordinate_gibbs_at_least_two}.
\end{proof}

\begin{remark}[Length-one boundary of the two-site construction]%
    \label{rem:mpdo_topological_gibbs_length_one_boundary}
    The local-interaction definition in \cite{Cirac2016MPDO_arXiv} lets a
    local operator act on the first two spins and translates it around the
    periodic chain. Thus its construction has domain $L\geq2$. Periodic
    identification of site labels supplies neither a map from two-site
    operators to one-site operators nor a multiplicity for the coincident
    edge.

    The source states its Gibbs-decomposition theorem without an explicit
    lower bound on $L$. Its $L=1$ instance requires an additional convention.
    The physical-space result here establishes the decomposition for
    $L\geq2$; no length-one statement is asserted.
    % See docs/paper-gaps/cpsv16_topological_gibbs_physical_complement.tex
    % and docs/paper-gaps/cpsv16_topological_gibbs_length_one.tex.
\end{remark}
